\section{Discussion and limitations}

\subsection{The formation advantage is a reach advantage}

$\Delta_\text{formation}$ is not a constant. It tracks how much of what a task needs
sits outside the roster, and it collapses when reach stops being scarce.

\begin{center}\small
\begin{tabular}{@{}lcc@{}}
\toprule
& $E = 0.3$ & $E = 0.7$ \\
\midrule
cold   & $+0.112$ & $+0.352$ \\
rework & $\mathbf{-0.037}$ & $+0.265$ \\
warm   & $+0.148$ & $+0.391$ \\
\bottomrule
\end{tabular}
\end{center}

At a low external fraction the gap is a third of its high-$E$ value, and on the rework
beat---where the task must return to a counterpart it has already used---the point
estimate reverses: the bounded network is ahead. Averaged over beats, the formation
effect halves on rework ($0.232 \to 0.114$).

We report the reversal as suggestive rather than established: it is a subcell point
estimate and we do not attach an interval to it. But the pattern across the table is
consistent and it bounds the headline. The open network wins because it can address
holders the bounded one cannot; where the bounded roster already contains the holders,
and where the work requires going back to them, that advantage is largely spent.

\subsection{A bounded roster as a wall, and what happens when it is not}

A bounded network in this design can address twenty holders and no others. That is a
truncated directory, not a private network: information reaches a closed group through
the people in it, one hop further out, slower and lossier. We therefore gave a roster
member the ability to forward one question to a contact of its own, minted under its own
authority rather than derived from the requester's, and re-ran.

The gap did not close. Across 64 episodes at sixteen per cell, relay moves the bounded
arm by $+0.042$ (95\% CI $[-0.137, 0.208]$), and the difference-in-differences against
the open arm---whose configuration is identical under both settings and which serves as
the negative control---is $+0.105$ (95\% CI $[-0.105, 0.322]$). The negative control
behaved: it moved $-0.063$, within 1.0 SE of nothing.

Two details are worth keeping. The same contrast read $+0.255$ at six episodes per cell
with an interval excluding zero and decayed monotonically as $n$ grew, so the small-$n$
estimate was noise. And excluding episodes that produced no artifact at all, the bounded
arm scores $0.734$ with relay against $0.554$ without: relay helps substantially
\emph{when it completes}, and costs enough exploration budget that seven episodes never
reached a build. Five of those were relay episodes and two were high first-contact-cost
episodes without it, so the failure belongs to any mechanism that consumes extra turns,
and what it exposes is a harness that answers budget exhaustion by refusing tools rather
than by forcing a build.

\subsection{Forming an edge is free here, and that is the main limitation}

Whether a link is worth its cost is the central question of strategic network formation
\cite{jackson1996}, where agents weigh the benefit of a connection against the price of
maintaining it. \textbf{We set that price to zero.} In this design, addressing a stranger
costs exactly what addressing a teammate costs.
The requester searches a directory and sends a message; there is no negotiation, no
context transfer, no format alignment, no establishing what a counterpart can actually
do.

This removes the economic reason relationships exist. An open arm with a hundred
addressable holders at zero acquisition cost is closer to an oracle than to a network,
and the formation effect we measure should be read as an upper bound: it is what open
discovery is worth \emph{when discovery is free}. A per-new-edge cost---an onboarding
exchange, a latency penalty, a token charge on first contact---is the single variable
most likely to move the ranking, and we have not run it. We flag it as the first
experiment a follow-up should do rather than as a caveat to be waved at.

\subsection{What the attribute null does and does not bound}

The null covers the weakest form of relational encoding: a sentence asserted in
context. It does not bound trust wired into routing, memory, retrieval, or permission
policy, which is where reputation lives in deployed systems and where it can change
what an agent is able to do rather than what it is told.

One feature of the design does work against detecting an attribute effect: no counterpart
lies strategically, so a reputation signal has no adversary to defend against, and a
property that matters mainly for defence cannot show its value where there is nothing to
defend against. The verification channel, by contrast, is present and used---90\% of beats
met a planted distractor and 94.5\% of encounters were resisted---so an assertion of prior
trust had a verification cost available to save and did not save it.

\subsection{Other limitations}

The cold-phase crossing follows from the fact allocation and is not a discovery; only its
location and cost are measured. Forced-value deliverables cannot capture whether an
artifact is well argued, which is the price of removing the judge. The kernel is a
reference implementation, not a deployment, and one model family carries the sweeps.
